\section*{The Propers: Themes and Movement}

\begin{massbox}{Receive Christ, lose the false life, live from his gift}
The Lectionary officially joins the Shunammite's welcome of Elisha to Jesus' promise that receiving his messenger receives him. Romans contributes an independent baptismal death-and-life strand, 1 Peter acclaims a called people brought into light, and the Missal frames these with adoption, truth, divinely effected mysteries, remembered benefits, ecclesial unity, and life in lasting charity.
\end{massbox}

\begin{center}
\textbf{Called into light $\longrightarrow$ Christ preferred $\longrightarrow$ messenger received $\longrightarrow$ life renewed}
\end{center}

\subsection*{1. Gift establishes the household\enspace\elementref{Entrance, Collect, Acclamation}}

The universal acclamation of Psalm 47 opens into the Collect's adoption by grace. The assembly does not generate light but is made children of light and asked to adhere to truth. First Peter gives this identity a vocation: those called from darkness proclaim God's mighty acts. Belonging becomes testimony.

\subsection*{2. Hospitality receives more than it controls\enspace\elementref{First Reading, Psalm, Gospel}}

The Shunammite recognizes Elisha, makes space, and receives a promise she did not manipulate. Matthew explicitly correlates receiving a prophet, righteous person, or little one with receiving Christ and the One who sent him. Psalm 89 answers in covenant mercy and blessed worship. Hospitality is receptive participation, never a transaction for guaranteed reward.

\subsection*{3. The cross orders every love\enspace\elementref{Gospel, Romans}}

Jesus' priority does not abolish family love; it judges and orders every creaturely attachment under love of God. Taking the cross and losing life are costly fidelity, not self-harm. Romans independently names the sacramental foundation of Christian death and life: Baptism joins believers to Christ's once-for-all Pasch and summons continuing death to sin.

\subsection*{4. God effects and Communion gives life\enspace\elementref{Prayer over Offerings, Communion, Prayer after Communion}}

The offering prayer attributes the mysteries' effect to God while asking that servants celebrate worthily. Communion may remember divine benefits or voice Christ's prayer for unity. The final prayer names the offered and received victim as life-giving and orders union toward lasting charity. This sacramental movement can illuminate Romans, but its Week XIII texts are shared across all cycles.

\begin{studybox}{Ordered love and prudent welcome}
Christ's priority never licenses neglect of children, spouse, parents, dependents, safeguarding, or lawful duties. The cross is not avoidable injury sought for itself. Hospitality does not suspend prudence or make access to vulnerable persons a religious entitlement. Baptismal identity grounds conversion; it does not make future sin impossible.
\end{studybox}
